\documentclass{beamer}

\usepackage{../packages/packages_mathalea}
\input{../competences/competences_latex.tex}
\usepackage{fontawesome}

\pgfplotsset{compat=1.18}

\begin{document}

\begin{frame}{\GetCompetence{Fonc16}}
Soit la parabole $y=\dfrac{3}{4}x^2-2x+1$.
\begin{itemize}
	\item Déterminer les coordonnées des points d'intersection de la parabole avec l'axe des abscisses.
	\item Déterminer les coordonnées des points d'intersection de la parabole avec l'axe des ordonnées.
	\item Déterminer les coordonnées du sommet de la parabole.
	\item Déterminer l'équation de l'axe de symétrie de la parabole.
	\item Déterminer l'écriture canonique de la parabole.
	\item Déterminer l'écriture factorisée de la parabole.
\end{itemize}

\onslide<2->{
\begin{multicols}{2}
\begin{itemize}
	\item $R_1\left(\dfrac{2}{3};0\right)$ et $R_2(2;0)$
	\item $(0;1)$
	\item $S\left(\dfrac{4}{3};-\dfrac{1}{3}\right)$
	\item $x=\dfrac{4}{3}$
	\item $y=\dfrac{3}{4}\left(x-\dfrac{4}{3}\right)^2-3$
	\item $y=\dfrac{3}{4}\left(x-\dfrac{2}{3}\right)(x-2)$
\end{itemize}
\end{multicols}
    }
\end{frame}
\end{document}
